\chapter{Tipos e Valores}

\hayashi{} é dinamicamente tipado: variáveis não têm tipo declarado, mas cada valor
carrega seu tipo em tempo de execução. O sistema de tipos é coerente — operações entre
tipos incompatíveis produzem erros claros em vez de coerções silenciosas.

\section{Tipos primitivos}

\subsection{Int}

Inteiro de 64 bits com sinal. Literais inteiros são escritos sem ponto decimal:

\begin{lstlisting}
let n = 42
let negativo = -7
let grande = 1_000_000    // _ como separador visual
\end{lstlisting}

\subsection{Float}

Ponto flutuante de 64 bits (IEEE 754). Requer ponto decimal ou notação científica:

\begin{lstlisting}
let pi = 3.14159
let e  = 2.71828
let pequeno = 1.5e-10
\end{lstlisting}

Operações aritméticas entre \texttt{Int} e \texttt{Float} promovem para \texttt{Float}:

\begin{lstlisting}
let x = 3 / 2      // 1  (divisão inteira)
let y = 3 / 2.0    // 1.5
let z = 3.0 / 2    // 1.5
\end{lstlisting}

Divisão de \texttt{Float} por zero retorna \texttt{inf} ou \texttt{-inf}, nunca
um erro — seguindo IEEE 754:

\begin{lstlisting}
display 1.0 / 0.0    // inf
display -1.0 / 0.0   // -inf
display 0.0 / 0.0    // nan
\end{lstlisting}

\begin{nota}
  Divisão de \texttt{Int} por zero produz erro de runtime. Divisão de \texttt{Float}
  por zero segue IEEE 754 e retorna \texttt{inf}.
\end{nota}

\subsection{Bool}

Verdadeiro ou falso. Literais: \lstinline{true} e \lstinline{false}.

\begin{lstlisting}
let ativo = true
let encerrado = false
\end{lstlisting}

Operadores lógicos: \lstinline{and}, \lstinline{or}, \lstinline{not}.

\subsection{Nil}

Ausência de valor. Algumas funções retornam \lstinline{nil} quando não há resultado
relevante (por exemplo, \lstinline{push} que muta uma lista no lugar).

\begin{lstlisting}
let nada = nil
display nada    // nil
\end{lstlisting}

\section{Strings}

Strings são sequências de caracteres Unicode delimitadas por aspas duplas:

\begin{lstlisting}
let nome = "Hayashi"
let caminho = "/home/usuario/dados.csv"
\end{lstlisting}

Sequências de escape: \texttt{\textbackslash n} (nova linha), \texttt{\textbackslash t}
(tab), \texttt{\textbackslash"} (aspas), \texttt{\textbackslash\textbackslash} (barra).

F-strings permitem interpolação de expressões:

\begin{lstlisting}
let ano = 2026
display f"Publicado em {ano}"   // Publicado em 2026
\end{lstlisting}

\section{Listas}

Coleções ordenadas de valores de qualquer tipo, delimitadas por colchetes:

\begin{lstlisting}
let numeros = [1, 2, 3, 4, 5]
let misto   = [1, "dois", true, nil]
let vazia   = []
\end{lstlisting}

\section{Dicionários}

Mapeamentos chave-valor. Chaves e valores podem ser de qualquer tipo:

\begin{lstlisting}
let config = {"host": "localhost", "porta": 5432}
let vazio  = {}
\end{lstlisting}

Acesso por chave:

\begin{lstlisting}
config["host"]      // "localhost"
\end{lstlisting}

\section{DataFrame}

O tipo central de \hayashi{}. Um DataFrame é uma tabela bidimensional com colunas
nomeadas. Cada coluna é um vetor de \texttt{Float} (valores ausentes representados
por \texttt{NaN}).

\begin{lstlisting}[caption={Criando um DataFrame inline}]
input x y
  1  2
  3  4
  5  6
end
\end{lstlisting}

DataFrames são discutidos em profundidade no Capítulo~\ref{cap:dataframes}.

\section{Conversões}

\hayashi{} não faz coerção implícita entre tipos. Conversões devem ser explícitas:

\begin{lstlisting}
let s = str(42)        // "42"
let n = int("7")       // 7
let f = float("3.14")  // 3.14
\end{lstlisting}

Tentar somar um \texttt{Int} com uma \texttt{String} produz erro de tipo:

\begin{lstlisting}
42 + "3"
// error: type mismatch — expected Int or Float, got String
\end{lstlisting}

\section{Verificação de Tipos}

Para verificar dinamicamente o tipo de um valor, \hayashi{} oferece funções embutidas de verificação que retornam um booleano:

\begin{lstlisting}
is_nil(nil)         // true
is_int(42)          // true
is_float(3.14)      // true
is_bool(true)       // true
is_str("ola")       // true  (ou is_string)
is_list([1, 2])     // true
is_dict({"a": 1})   // true
is_df(df)           // true  (ou is_dataframe)
is_fn(|x| x)        // true  (ou is_function)
\end{lstlisting}

Além disso, a função \lstinline{type(x)} (ou \lstinline{typeof(x)}) retorna o tipo do valor como uma string (ex: \lstinline{"int"}, \lstinline{"float"}, \lstinline{"nil"}, etc.).
